\documentclass[10pt,a4paper]{amsart}
\usepackage[margin=19mm]{geometry}
\usepackage{amsmath,amssymb,mathtools}
\usepackage[scr=rsfs,cal=euler]{mathalfa}
\usepackage{xfrac}
\usepackage[unicode,hidelinks,bookmarks=false]{hyperref}
\newcommand{\ie}{i.e.\ }
\newcommand{\cf}{cf.\ }
\newcommand{\resp}{respectively\ }
\newcommand{\Subsection}{\subsection}
\providecommand{\nobreakdash}{\nobreak\hbox{-}}
\setlength{\emergencystretch}{2em}
\multlinegap0pt
\usepackage{enumerate}
\usepackage{amssymb}
\usepackage{mathtools}
\usepackage{float}
\usepackage{xcolor}
\newtheorem{lemma}{Lemma}
\newtheorem{theorem}{Theorem}
\usepackage{cleveref}
\crefname{lemma}{Lemma}{Lemmas}
\crefname{theorem}{Theorem}{Theorems}
\newcommand{\NN}{\mathbb Z_{\geq0}}
\newcommand{\RR}{\mathbb R}
\title{Real-Rootedness and Gamma-Positivity for a Variation of the Morris Constant Term}
\author{Feihu Liu$^{\color{blue} \dag}$ and Zihao Zhang$^{\color{blue} \S}$}
\date{Source: arXiv:2609.15201v1 (CC BY 4.0)}
\begin{document}
\maketitle

\noindent\emph{Source and licence.} Real-Rootedness and Gamma-Positivity for a Variation of the Morris Constant Term. Version arXiv:2609.15201v1 (CC BY 4.0), distributed by arXiv under the Creative Commons Attribution 4.0 International licence, http://creativecommons.org/licenses/by/4.0/, as recorded in the arXiv licence metadata for this exact version and shown on the abstract page https://arxiv.org/abs/2609.15201v1. Full licence notice: \texttt{licenses/arxiv-2609.15201-CC-BY-4.0.txt}.\par\smallskip

\noindent\textit{Editorial scope.} Counted unit: the article's theorem st:stable-eigenfunction (source0.tex lines 1559--1574). The article's complete original proof of it is delivered with it (source0.tex lines 1575--1676, one \texttt{proof} environment of 102 lines). No mathematics is written, repaired or re-derived: the statement and the proof are the article's own lines, in the article's own notation, and no retained line is substituted or re-flowed.  Retained uncounted as the source-local context the proof needs: lines 997--1001; lines 1082--1091; lines 1387--1390; lines 1457--1468; lines 1514--1521; lines 1516--1518.  Nothing was omitted from any retained span, so the package carries no omission record.  No retained line contains a citation, so the wrapper carries no bibliography and no bibliography entry.  No retained line contains a figure environment, an image-inclusion command or drawing code, so no image is delivered and the package declares no asset dependency.  Editorial formatting only, in addition: an AMS \texttt{amsart} preamble that restates the article's own declarations and macros used by the retained lines, namely the environments \texttt{lemma}, \texttt{theorem} on separate counters, together with the standard packages those lines need.  The retained environments print in this document as Lemma 1, Theorem 1 in order of appearance, whereas the article prints the same items with the numbering of its own full document.  The complete native source of the article as distributed in the arXiv e-print for this exact version is bundled unchanged as \texttt{sources/210398/source0.tex}.
\begin{lemma}\label{st:limit-closure}
Let $F_j$ be stable polynomials in a fixed number of variables, of
uniformly bounded total degree.  If they converge coefficientwise to
$F$, then $F$ is either stable or identically zero.
\end{lemma}

\begin{lemma}\label{st:elementary-unitriangular}
The polynomial $E_\lambda$ is real stable, and
\begin{equation}\label{st:elementary-triangular}
E_\lambda=m_\lambda+
 \sum_{\mu\prec\lambda}a_{\lambda\mu}m_\mu,
\qquad a_{\lambda\mu}\in\NN.
\end{equation}
The polynomials $E_\mu$ with $\mu\preceq\lambda$ therefore form
a basis of $\mathcal V_\lambda$.
\end{lemma}

\begin{theorem}\label{st:evolution-preserves}
For $t\geq0$ and $R\in\NN$, the operator $e^{t\mathcal D_q}$
on $\mathcal S_{q,\leq R}$ preserves stability and real stability.
\end{theorem}

\begin{lemma}\label{st:operator-triangular}
For every partition $\nu$ of length at most $q$, we have 
\begin{equation}\label{st:triangular-operator}
\mathcal D_q m_\nu=\varepsilon_\nu m_\nu+\sum_{\mu\prec\nu}c_{\nu\mu}m_\mu,
\end{equation}
where $c_{\nu\mu}\in\RR$ and
\begin{align}
\varepsilon_\nu&=\sum_{i=1}^q\nu_i\bigl(2(q-i)-\nu_i+1\bigr)\label{st:eigenvalue}\\
&=2\sum_{(i,j)\in[\nu]}(q+1-i-j).\label{st:eigenvalue-diagram}
\end{align}
In particular, $\mathcal V_\lambda$ is invariant under $\mathcal D_q$.
\end{lemma}

\begin{lemma}\label{st:spectral-gap}
Suppose that $\lambda$ has length at most $q$ and satisfies
\begin{equation}\label{st:adjacent-drop}
\lambda_i-\lambda_{i+1}\leq1\qquad(1\leq i<q).
\end{equation}
Then $\varepsilon_\lambda>\varepsilon_\mu$ for every
$\mu\prec\lambda$ of length at most $q$.
\end{lemma}

\begin{equation}
\lambda_i-\lambda_{i+1}\leq1\qquad(1\leq i<q).
\end{equation}

\begin{theorem}\label{st:stable-eigenfunction}
Let $q\geq1$ and let $\lambda$ have length at most $q$ and satisfy
\eqref{st:adjacent-drop}.  There is a unique polynomial
$P_\lambda\in\mathcal V_\lambda$ of the form
\begin{equation}\label{st:monic-eigenfunction-form}
P_\lambda=m_\lambda+\sum_{\mu\prec\lambda}p_{\lambda\mu}m_\mu
\end{equation}
satisfying
\begin{equation}\label{st:monic-eigenfunction-equation}
\mathcal D_qP_\lambda=\varepsilon_\lambda P_\lambda.
\end{equation}
It has real coefficients and is real stable.  Consequently, every
polynomial in $\mathcal V_\lambda$ with a nonzero $m_\lambda$
coefficient and satisfying \eqref{st:monic-eigenfunction-equation}
is stable.
\end{theorem}

\begin{proof}
Order the finite dominance downset with larger partitions first,
extending dominance to a total order, and set
$p_{\lambda\lambda}=1$.
By \cref{st:operator-triangular}, we have
$$\mathcal D_q m_\nu=\varepsilon_\nu m_\nu+\sum_{\mu\prec\nu}c_{\nu\mu}m_\mu$$
According to
$$\mathcal D_q\left(m_\lambda+\sum_{\mu\prec\lambda}p_{\lambda\mu}m_\mu\right)=\varepsilon_\lambda \left(m_\lambda+\sum_{\mu\prec\lambda}p_{\lambda\mu}m_\mu\right),$$
the $m_\mu$ coefficient of the eigenvalue equation is
\[\varepsilon_\mu p_{\lambda\mu}+\sum_{\mu\prec\nu\preceq\lambda}c_{\nu\mu}p_{\lambda\nu}=\varepsilon_\lambda p_{\lambda\mu}.
\]
For $\mu=\lambda$ this is automatic; otherwise it gives
\begin{equation}\label{st:eigenfunction-recursion}
p_{\lambda\mu}
=\frac{\displaystyle\sum_{\mu\prec\nu\preceq\lambda}
        c_{\nu\mu}p_{\lambda\nu}}
       {\varepsilon_\lambda-\varepsilon_\mu}.
\end{equation}
All coefficients on the right occur earlier in the order, and every
denominator is positive by \cref{st:spectral-gap}.
Thus the recursion gives existence, uniqueness, and real coefficients.

To prove stability, set
\begin{equation}\label{st:normalized-evolution}
P_t=e^{-t\varepsilon_\lambda}
       e^{t\mathcal D_q}E_\lambda,
\qquad t\geq0.
\end{equation}
By \cref{st:elementary-unitriangular} and \cref{st:evolution-preserves}, each $P_t$ is real stable.
The invariance of $\mathcal V_\lambda$ keeps $P_t$ in that space,
and differentiation gives
\begin{equation}\label{st:normalized-differential-equation}
P_t'=(\mathcal D_q-\varepsilon_\lambda I)P_t.
\end{equation}
Triangularity shows that the $m_\lambda$ coefficient remains
at its initial value $1$.  Write
\[
P_t=\sum_{\mu\preceq\lambda}u_\mu(t)m_\mu,
\qquad u_\lambda(t)=1.
\]
For $\mu\prec\lambda$, coefficient comparison according to \eqref{st:normalized-differential-equation} gives
\begin{equation}\label{st:coefficient-evolution}
u_\mu'(t)
=-(\varepsilon_\lambda-\varepsilon_\mu)u_\mu(t)
  +\sum_{\mu\prec\nu\preceq\lambda}
      c_{\nu\mu}u_\nu(t).
\end{equation}
Here $c_{\nu\mu}=0$ for absent terms.

We prove coefficient convergence in the same total order.  Suppose that all preceding coefficients of $u_{\mu}(t)$ converge.  Then
$v_\mu(t)=\sum_{\mu\prec\nu\preceq\lambda}c_{\nu\mu}u_\nu(t)$
is bounded and has a limit $v_\mu(\infty)$.
Set $\delta_\mu=\varepsilon_\lambda-\varepsilon_\mu>0$.
According to \eqref{st:coefficient-evolution}, we have
$$u'_\mu(t) + \delta_\mu u_\mu(t) = v_\mu(t)$$
Thus, we have $e^{\delta_\mu t} u'_\mu(t) + e^{\delta_\mu t}\delta_\mu u_\mu(t) =e^{\delta_\mu t} v_\mu(t)$.
Therefore, we obtain
$$(e^{\delta_\mu t} u_\mu(t))'=e^{\delta_\mu t} v_\mu(t).$$
Integrating from $0$ to $t$ and then dividing by $e^{\delta_\mu t}$, we obtain 
\begin{equation}\label{st:variation-of-constants}
u_\mu(t)=e^{-\delta_\mu t}u_\mu(0)
 +\int_0^t e^{-\delta_\mu(t-s)}v_\mu(s)\,\mathrm ds.
\end{equation}
It follows that
\begin{align*}
u_\mu(t)-\frac{v_\mu(\infty)}{\delta_\mu}=e^{-\delta_\mu t}\left(u_\mu(0)-\frac{v_\mu(\infty)}{\delta_\mu}\right)+\int_0^t e^{-\delta_\mu(t-s)}(v_\mu(s)-v_\mu(\infty))\,\mathrm ds.
\end{align*}
For $T>0$, set
$K_T=\sup_{0\leq s\leq T}|v_\mu(s)-v_\mu(\infty)|$.
When $t\geq T$, the part of the last integral over $[0,T]$ is
bounded in modulus by $TK_Te^{-\delta_\mu(t-T)}$, and the
remaining part by
\[
\sup_{s\geq T}|v_\mu(s)-v_\mu(\infty)|
 \int_T^t e^{-\delta_\mu(t-s)}\,\mathrm ds
\leq\frac1{\delta_\mu}
 \sup_{s\geq T}|v_\mu(s)-v_\mu(\infty)|.
\]
First let $t\to\infty$ and then $T\to\infty$.
Both bounds vanish, proving
$u_\mu(t)\to v_\mu(\infty)/\delta_\mu$ and completing the induction.
Thus $u_\mu(t)$ converges. This proves convergence of all coefficients.

Writing $\ell_\mu=\lim_{t\to\infty}u_\mu(t)$, we obtain
$\ell_\lambda=1$ and
\begin{equation}\label{st:limit-recursion}
\ell_\mu
=\frac{v_\mu(\infty)}{\delta_\mu}
=\frac{\displaystyle\sum_{\mu\prec\nu\preceq\lambda}
                 c_{\nu\mu}\ell_\nu}
       {\varepsilon_\lambda-\varepsilon_\mu}
\qquad(\mu\prec\lambda).
\end{equation}
This is exactly \eqref{st:eigenfunction-recursion}; uniqueness
identifies the limit with $P_\lambda$.  The limit is nonzero,
since its $m_\lambda$ coefficient is $1$. \cref{st:limit-closure}, applied to $P_1,P_2,\ldots$,
therefore proves real stability.

Finally, dividing any polynomial in the last assertion by its
nonzero $m_\lambda$ coefficient gives $P_\lambda$ by uniqueness.
A nonzero scalar multiple has the same zero set.
\end{proof}

\end{document}
